% =====================================================================
% Abbildung: Verbindungsmodell mit Ports "part def PerceptionSystem"
% Passend zu Kapitel 9.4/9.6 (Listings "port def ...", "part def
% PerceptionSystem { ... connection c1 ... connection c2 ... }")
%
% Einbindung im Buch, z. B. in chapters/09_ibd.tex:
%
%   \begin{figure}[H]
%     \centering
%     % =====================================================================
% Abbildung: Verbindungsmodell mit Ports "part def PerceptionSystem"
% Passend zu Kapitel 9.4/9.6 (Listings "port def ...", "part def
% PerceptionSystem { ... connection c1 ... connection c2 ... }")
%
% Einbindung im Buch, z. B. in chapters/09_ibd.tex:
%
%   \begin{figure}[H]
%     \centering
%     % =====================================================================
% Abbildung: Verbindungsmodell mit Ports "part def PerceptionSystem"
% Passend zu Kapitel 9.4/9.6 (Listings "port def ...", "part def
% PerceptionSystem { ... connection c1 ... connection c2 ... }")
%
% Einbindung im Buch, z. B. in chapters/09_ibd.tex:
%
%   \begin{figure}[H]
%     \centering
%     % =====================================================================
% Abbildung: Verbindungsmodell mit Ports "part def PerceptionSystem"
% Passend zu Kapitel 9.4/9.6 (Listings "port def ...", "part def
% PerceptionSystem { ... connection c1 ... connection c2 ... }")
%
% Einbindung im Buch, z. B. in chapters/09_ibd.tex:
%
%   \begin{figure}[H]
%     \centering
%     \input{figures/fig_ports_connections}
%     \caption{Verbindungsmodell des \texttt{PerceptionSystem}. Die kleinen
%       Quadrate an den Boxraendern sind Ports. Das Tilde-Zeichen markiert
%       einen konjugierten Port, der empfaengt statt sendet.}
%     \label{fig:ports-connections}
%   \end{figure}
%
% Benoetigte Pakete (im Buch bereits in main.tex geladen): tikz
% Benoetigte TikZ-Bibliotheken: positioning, arrows.meta, shadows, calc
% =====================================================================

\begin{tikzpicture}[
    every node/.style={font=\small},
    part/.style={
        rectangle, rounded corners=1.5pt, draw=black!70, thick,
        minimum width=30mm, minimum height=12mm,
        align=center, fill=blue!8, drop shadow
    },
    port/.style={rectangle, draw=black, fill=black!75, minimum width=2.2mm,
        minimum height=2.2mm, inner sep=0pt},
    portlbl/.style={font=\tiny},
    flow/.style={-{Stealth[length=2.2mm]}, thick, orange!70!black}
]

% --- Systemteile ---
\node[part] (camera) at (0, 3)   {\texttt{camera :}\\\texttt{Camera}};
\node[part] (lidar)  at (0, 0)   {\texttt{lidar :}\\\texttt{Lidar}};
\node[part] (detect) at (7, 3)   {\texttt{objectDetector :}\\\texttt{ObjectDetector}};
\node[part] (obst)   at (7, 0)   {\texttt{obstacleDetector :}\\\texttt{ObstacleDetector}};

% --- Ports (kleine Quadrate am Rand) ---
\node[port] (p-cam-out)   at ($(camera.east)$)   {};
\node[port] (p-lidar-out) at ($(lidar.east)$)    {};
\node[port] (p-det-in)    at ($(detect.west)$)   {};
\node[port] (p-det-out)   at ($(detect.east)$)   {};
\node[port] (p-obst-in)   at ($(obst.west)$)     {};

\node[portlbl, above=0.5mm of p-cam-out]   {\texttt{imageOut}};
\node[portlbl, above=0.5mm of p-lidar-out] {\texttt{scanOut}};
\node[portlbl, above=0.5mm of p-det-in]    {\texttt{imageIn} \textcolor{orange!70!black}{$\sim$}};
\node[portlbl, above=0.5mm of p-det-out]   {\texttt{objectsOut}};
\node[portlbl, above=0.5mm of p-obst-in]   {\texttt{scanIn} \textcolor{orange!70!black}{$\sim$}};

% --- Verbindungen (connection c1 / c2) ---
\draw[flow] (p-cam-out)   -- node[above, midway, font=\tiny]{c1} (p-det-in);
\draw[flow] (p-lidar-out) -- node[above, midway, font=\tiny]{c2} (p-obst-in);

% --- Weiterleitung an Navigation (objectsOut, im Listing nicht verbunden) ---
\draw[flow, dashed] (p-det-out) -- ++(2.2,0)
    node[right, font=\tiny, align=left] {an\\Navigation};

% --- Legende ---
\node[port, below=16mm of lidar, xshift=-4mm] (leg1) {};
\node[portlbl, right=1mm of leg1] {Port};
\node[portlbl, below=2mm of leg1, xshift=6mm] (legtilde) {\textcolor{orange!70!black}{$\sim$} = konjugierter Port (empfaengt statt sendet)};
\draw[flow] ([yshift=-8mm]leg1.south) -- ++(6mm,0)
    node[right=1mm, font=\tiny] {Verbindung (\texttt{connection})};

\end{tikzpicture}

%     \caption{Verbindungsmodell des \texttt{PerceptionSystem}. Die kleinen
%       Quadrate an den Boxraendern sind Ports. Das Tilde-Zeichen markiert
%       einen konjugierten Port, der empfaengt statt sendet.}
%     \label{fig:ports-connections}
%   \end{figure}
%
% Benoetigte Pakete (im Buch bereits in main.tex geladen): tikz
% Benoetigte TikZ-Bibliotheken: positioning, arrows.meta, shadows, calc
% =====================================================================

\begin{tikzpicture}[
    every node/.style={font=\small},
    part/.style={
        rectangle, rounded corners=1.5pt, draw=black!70, thick,
        minimum width=30mm, minimum height=12mm,
        align=center, fill=blue!8, drop shadow
    },
    port/.style={rectangle, draw=black, fill=black!75, minimum width=2.2mm,
        minimum height=2.2mm, inner sep=0pt},
    portlbl/.style={font=\tiny},
    flow/.style={-{Stealth[length=2.2mm]}, thick, orange!70!black}
]

% --- Systemteile ---
\node[part] (camera) at (0, 3)   {\texttt{camera :}\\\texttt{Camera}};
\node[part] (lidar)  at (0, 0)   {\texttt{lidar :}\\\texttt{Lidar}};
\node[part] (detect) at (7, 3)   {\texttt{objectDetector :}\\\texttt{ObjectDetector}};
\node[part] (obst)   at (7, 0)   {\texttt{obstacleDetector :}\\\texttt{ObstacleDetector}};

% --- Ports (kleine Quadrate am Rand) ---
\node[port] (p-cam-out)   at ($(camera.east)$)   {};
\node[port] (p-lidar-out) at ($(lidar.east)$)    {};
\node[port] (p-det-in)    at ($(detect.west)$)   {};
\node[port] (p-det-out)   at ($(detect.east)$)   {};
\node[port] (p-obst-in)   at ($(obst.west)$)     {};

\node[portlbl, above=0.5mm of p-cam-out]   {\texttt{imageOut}};
\node[portlbl, above=0.5mm of p-lidar-out] {\texttt{scanOut}};
\node[portlbl, above=0.5mm of p-det-in]    {\texttt{imageIn} \textcolor{orange!70!black}{$\sim$}};
\node[portlbl, above=0.5mm of p-det-out]   {\texttt{objectsOut}};
\node[portlbl, above=0.5mm of p-obst-in]   {\texttt{scanIn} \textcolor{orange!70!black}{$\sim$}};

% --- Verbindungen (connection c1 / c2) ---
\draw[flow] (p-cam-out)   -- node[above, midway, font=\tiny]{c1} (p-det-in);
\draw[flow] (p-lidar-out) -- node[above, midway, font=\tiny]{c2} (p-obst-in);

% --- Weiterleitung an Navigation (objectsOut, im Listing nicht verbunden) ---
\draw[flow, dashed] (p-det-out) -- ++(2.2,0)
    node[right, font=\tiny, align=left] {an\\Navigation};

% --- Legende ---
\node[port, below=16mm of lidar, xshift=-4mm] (leg1) {};
\node[portlbl, right=1mm of leg1] {Port};
\node[portlbl, below=2mm of leg1, xshift=6mm] (legtilde) {\textcolor{orange!70!black}{$\sim$} = konjugierter Port (empfaengt statt sendet)};
\draw[flow] ([yshift=-8mm]leg1.south) -- ++(6mm,0)
    node[right=1mm, font=\tiny] {Verbindung (\texttt{connection})};

\end{tikzpicture}

%     \caption{Verbindungsmodell des \texttt{PerceptionSystem}. Die kleinen
%       Quadrate an den Boxraendern sind Ports. Das Tilde-Zeichen markiert
%       einen konjugierten Port, der empfaengt statt sendet.}
%     \label{fig:ports-connections}
%   \end{figure}
%
% Benoetigte Pakete (im Buch bereits in main.tex geladen): tikz
% Benoetigte TikZ-Bibliotheken: positioning, arrows.meta, shadows, calc
% =====================================================================

\begin{tikzpicture}[
    every node/.style={font=\small},
    part/.style={
        rectangle, rounded corners=1.5pt, draw=black!70, thick,
        minimum width=30mm, minimum height=12mm,
        align=center, fill=blue!8, drop shadow
    },
    port/.style={rectangle, draw=black, fill=black!75, minimum width=2.2mm,
        minimum height=2.2mm, inner sep=0pt},
    portlbl/.style={font=\tiny},
    flow/.style={-{Stealth[length=2.2mm]}, thick, orange!70!black}
]

% --- Systemteile ---
\node[part] (camera) at (0, 3)   {\texttt{camera :}\\\texttt{Camera}};
\node[part] (lidar)  at (0, 0)   {\texttt{lidar :}\\\texttt{Lidar}};
\node[part] (detect) at (7, 3)   {\texttt{objectDetector :}\\\texttt{ObjectDetector}};
\node[part] (obst)   at (7, 0)   {\texttt{obstacleDetector :}\\\texttt{ObstacleDetector}};

% --- Ports (kleine Quadrate am Rand) ---
\node[port] (p-cam-out)   at ($(camera.east)$)   {};
\node[port] (p-lidar-out) at ($(lidar.east)$)    {};
\node[port] (p-det-in)    at ($(detect.west)$)   {};
\node[port] (p-det-out)   at ($(detect.east)$)   {};
\node[port] (p-obst-in)   at ($(obst.west)$)     {};

\node[portlbl, above=0.5mm of p-cam-out]   {\texttt{imageOut}};
\node[portlbl, above=0.5mm of p-lidar-out] {\texttt{scanOut}};
\node[portlbl, above=0.5mm of p-det-in]    {\texttt{imageIn} \textcolor{orange!70!black}{$\sim$}};
\node[portlbl, above=0.5mm of p-det-out]   {\texttt{objectsOut}};
\node[portlbl, above=0.5mm of p-obst-in]   {\texttt{scanIn} \textcolor{orange!70!black}{$\sim$}};

% --- Verbindungen (connection c1 / c2) ---
\draw[flow] (p-cam-out)   -- node[above, midway, font=\tiny]{c1} (p-det-in);
\draw[flow] (p-lidar-out) -- node[above, midway, font=\tiny]{c2} (p-obst-in);

% --- Weiterleitung an Navigation (objectsOut, im Listing nicht verbunden) ---
\draw[flow, dashed] (p-det-out) -- ++(2.2,0)
    node[right, font=\tiny, align=left] {an\\Navigation};

% --- Legende ---
\node[port, below=16mm of lidar, xshift=-4mm] (leg1) {};
\node[portlbl, right=1mm of leg1] {Port};
\node[portlbl, below=2mm of leg1, xshift=6mm] (legtilde) {\textcolor{orange!70!black}{$\sim$} = konjugierter Port (empfaengt statt sendet)};
\draw[flow] ([yshift=-8mm]leg1.south) -- ++(6mm,0)
    node[right=1mm, font=\tiny] {Verbindung (\texttt{connection})};

\end{tikzpicture}

%     \caption{Verbindungsmodell des \texttt{PerceptionSystem}. Die kleinen
%       Quadrate an den Boxraendern sind Ports. Das Tilde-Zeichen markiert
%       einen konjugierten Port, der empfaengt statt sendet.}
%     \label{fig:ports-connections}
%   \end{figure}
%
% Benoetigte Pakete (im Buch bereits in main.tex geladen): tikz
% Benoetigte TikZ-Bibliotheken: positioning, arrows.meta, shadows, calc
% =====================================================================

\begin{tikzpicture}[
    every node/.style={font=\small},
    part/.style={
        rectangle, rounded corners=1.5pt, draw=black!70, thick,
        minimum width=30mm, minimum height=12mm,
        align=center, fill=blue!8, drop shadow
    },
    port/.style={rectangle, draw=black, fill=black!75, minimum width=2.2mm,
        minimum height=2.2mm, inner sep=0pt},
    portlbl/.style={font=\tiny},
    flow/.style={-{Stealth[length=2.2mm]}, thick, orange!70!black}
]

% --- Systemteile ---
\node[part] (camera) at (0, 3)   {\texttt{camera :}\\\texttt{Camera}};
\node[part] (lidar)  at (0, 0)   {\texttt{lidar :}\\\texttt{Lidar}};
\node[part] (detect) at (7, 3)   {\texttt{objectDetector :}\\\texttt{ObjectDetector}};
\node[part] (obst)   at (7, 0)   {\texttt{obstacleDetector :}\\\texttt{ObstacleDetector}};

% --- Ports (kleine Quadrate am Rand) ---
\node[port] (p-cam-out)   at ($(camera.east)$)   {};
\node[port] (p-lidar-out) at ($(lidar.east)$)    {};
\node[port] (p-det-in)    at ($(detect.west)$)   {};
\node[port] (p-det-out)   at ($(detect.east)$)   {};
\node[port] (p-obst-in)   at ($(obst.west)$)     {};

\node[portlbl, above=0.5mm of p-cam-out]   {\texttt{imageOut}};
\node[portlbl, above=0.5mm of p-lidar-out] {\texttt{scanOut}};
\node[portlbl, above=0.5mm of p-det-in]    {\texttt{imageIn} \textcolor{orange!70!black}{$\sim$}};
\node[portlbl, above=0.5mm of p-det-out]   {\texttt{objectsOut}};
\node[portlbl, above=0.5mm of p-obst-in]   {\texttt{scanIn} \textcolor{orange!70!black}{$\sim$}};

% --- Verbindungen (connection c1 / c2) ---
\draw[flow] (p-cam-out)   -- node[above, midway, font=\tiny]{c1} (p-det-in);
\draw[flow] (p-lidar-out) -- node[above, midway, font=\tiny]{c2} (p-obst-in);

% --- Weiterleitung an Navigation (objectsOut, im Listing nicht verbunden) ---
\draw[flow, dashed] (p-det-out) -- ++(2.2,0)
    node[right, font=\tiny, align=left] {an\\Navigation};

% --- Legende ---
\node[port, below=16mm of lidar, xshift=-4mm] (leg1) {};
\node[portlbl, right=1mm of leg1] {Port};
\node[portlbl, below=2mm of leg1, xshift=6mm] (legtilde) {\textcolor{orange!70!black}{$\sim$} = konjugierter Port (empfaengt statt sendet)};
\draw[flow] ([yshift=-8mm]leg1.south) -- ++(6mm,0)
    node[right=1mm, font=\tiny] {Verbindung (\texttt{connection})};

\end{tikzpicture}
